\section{Coordinates and collision estimates}\label{sec:framework}

We first fix the normalization used in the estimates and explain the
singular velocity weight. Throughout the construction we rotate space
and velocity by the same orthogonal map, taking the original direction
\((1,\sqrt2,0)/\sqrt3\) to the third coordinate axis. The period lattice
becomes a rotation \(\Lambda\) of \(\mathbb Z^3\), still of covolume one.
The equation and all the properties in Definition~\ref{def:admissible}
are invariant under this change of coordinates. We work on
\(\R^3/\Lambda\), or in its periodic lift, until the last step of the
proof. Write \(x=(y,z)\in\R^2\times\R\) and
\(v=(\bar v,v_3)\). Dilations in space are always taken in the lift.

Let \(P\) denote the bilinear gain
\[
 P(f,g)(v)=\int_{\R^3\times S^2}
 |(v-v_*)\cdot\omega|f(v')g(v_*')\dd v_*\dd\omega.
\]
For the hard-sphere kernel it is symmetric in \(f,g\), as verified in
Lemma~\ref{lem:carleman}; in particular \(P(F,F)=Q^+(F,F)\).
The loss frequency is
\begin{equation}\label{eq:frequency}
 \nu_f(v)=2\pi\int_{\R^3}|v-p|f(p)\dd p,
 \qquad Q^-(F,F)=F\nu_F.
\end{equation}
A superscript \(b\), where \(b\ge1\), means that the collision kernel is
multiplied by
\begin{equation}\label{eq:cutoff}
 \min\{1,b/|v-v_*|\}.
\end{equation}
Its value at zero is immaterial. The same factor is used in gain and
loss; thus \(\nu_f^b(v)=2\pi\int\min\{b,|v-p|\}f(p)\dd p\).
It preserves particle interchange and the elastic change of variables,
and is exactly one on any fixed bounded velocity region for all
sufficiently large \(b\).

Fix \(1<q<2\) and \(m=100\). Our velocity weight is
\begin{equation}\label{eq:weight}
 H(v)=e^{-|v|^2}\langle v\rangle^{-m}W(v),
 \qquad W(v)=1+|\bar v|^{-q},\qquad
 \|f\|_H=\esssup_{x,v}\frac{|f(x,v)|}{H(v)}.
\end{equation}
The value assigned on the velocity axis is irrelevant. All supremum
estimates are essential estimates away from that axis. Since its
transverse dimension is two and \(q<2\), \(H\) is integrable with every
polynomial velocity weight; it also belongs to \(L^{1+\varepsilon}\)
for small enough \(\varepsilon>0\). The singularity permits gain
estimates uniform as a cold velocity distribution collapses onto the
axis. The inequality \(q>1\) is used in the corresponding plane
integrals.

Constants denoted \(C\) may depend on fixed construction parameters,
but not on the approximation index or the shrinking horizon unless
explicitly stated. The choices of the four parameters governing
branching growth will be made in Section~\ref{sec:branching}; all earlier
arguments hold for any fixed parameters in the indicated ranges. The
remaining small parameters are then chosen in the order specified by
the continuation proof.

\paragraph{Transport representatives.}
An integrable distributional equation \(\D f=R\) is read in free
coordinates, \(f^\sharp(t,x,v)=f(t,x+tv,v)\). Its time sections are
absolutely continuous for almost every \((x,v)\), with derivative
\(R^\sharp\). We use these representatives in path inequalities.
Fubini's theorem makes such inequalities valid for almost every
endpoint at any prescribed time. At the finite approximations the
gain kernels have absolutely continuous hot-input marginals; together
with free translations this permits iteration on common full-measure
endpoint sets. The later limiting path formulas will be established
from their integrable sources. No assertion depends on values assigned
to the velocity axis or other null sets.
